\chapter*{}
\noindent{\textbf{RESUMO}}

\noindent{Este trabalho apresenta o desenvolvimento do GUY, uma extensão para o Visual Studio Code que gera automaticamente o Grafo de Fluxo de Controle (GFC) e calcula a Complexidade Ciclomática. A pesquisa compreendeu a análise de ferramentas correlatas, a definição dos requisitos, a prototipagem e a implementação da extensão em TypeScript, com Tree-sitter e Webview. O protótipo foi verificado com três exemplos controlados em Python. Neles, foram obtidos, respectivamente, \(V(G)=4\), \(V(G)=2\) e \(V(G)=3\), com a mesma quantidade de caminhos independentes em cada caso. Para as estruturas avaliadas, os resultados apresentaram coerência entre o código de entrada, a topologia dos grafos e a fórmula da métrica. Conclui-se que o protótipo reuniu os recursos definidos para geração e visualização interativa do GFC, cálculo de \(V(G)\) e navegação entre grafo e código-fonte no ambiente local do desenvolvedor. A avaliação, por estar restrita a três exemplos, constitui uma evidência inicial de funcionamento e não permite generalizar os resultados para outros programas ou contextos de uso.}

\vspace{\onelineskip}

% todas em letras minúsculas, separadas por ponto e vírgula (;)
\noindent{\textbf{Palavras-chave}: grafo de fluxo de controle; complexidade ciclomática; teste estrutural; visual studio code.}

\vspace{\onelineskip}

\noindent{\textbf{ABSTRACT}}

\noindent{This work presents the development of GUY, a Visual Studio Code extension that automatically generates Control Flow Graphs (CFGs) and calculates Cyclomatic Complexity. The research comprised an analysis of related tools, requirements definition, prototyping, and implementation of the extension in TypeScript using Tree-sitter and Webview. The prototype was examined using three controlled Python examples. They produced, respectively, \(V(G)=4\), \(V(G)=2\), and \(V(G)=3\), with the same number of independent paths in each case. For the evaluated structures, the results were consistent with the input code, graph topology, and metric formula. The prototype therefore incorporated the intended features for interactive CFG generation and visualization, \(V(G)\) calculation, and navigation between the graph and source code in the developer's local environment. Because the evaluation was limited to three examples, it provides initial evidence of operation and does not support generalization to other programs or usage contexts.}

\vspace{\onelineskip}

\noindent{\textbf{Keywords}: control flow graph; cyclomatic complexity; structural testing; visual studio code.}
